\documentclass[11pt,a4paper]{article}

% ---------------------------------------------------------
% PREAMBLE (stable, warning-free, Overleaf-safe)
% ---------------------------------------------------------
\usepackage[utf8]{inputenc}
\usepackage[T1]{fontenc}
\usepackage{lmodern}
\usepackage{amsmath, amssymb, amsfonts}
\usepackage{bm}
\usepackage{geometry}
\usepackage{graphicx}
\usepackage{hyperref}
\usepackage{cite}
\usepackage{authblk}
\usepackage{physics}
\usepackage{mathtools}

\geometry{margin=2.4cm}

\title{\textbf{Companion XLIII: The Velocity Divergence Field in the Relaxation-Driven Cyclic Cosmology}}
\author{Michael Lehmann \\ \texttt{mi.lehmann@gmx.de}}
\date{April 2026}

\begin{document}
\maketitle

\begin{abstract}
The velocity divergence field encodes the dynamical response of matter to the 
gravitational potential and provides a direct probe of the growth of structure. 
In the standard cosmological model, the divergence evolves according to the linear 
growth rate and the continuity equation, producing a velocity field whose coherence 
and amplitude reflect the underlying matter distribution. In the Relaxation-Driven 
Cyclic Cosmology (RDCC), the scale-dependent evolution of the gravitational 
potential modifies the growth of the velocity divergence, the amplitude of bulk 
flows, and the coupling between density and velocity fields. This Companion 
develops a continuous description of the RDCC velocity divergence field, deriving 
how the relaxation scale influences the growth rate, the coherence of large-scale 
flows, and the dynamical architecture of the cosmic web. The resulting framework 
provides a new dynamical perspective on RDCC structure formation.
\end{abstract}
% ---------------------------------------------------------
\section{Introduction}

The cosmic velocity field is a dynamical tracer of the gravitational potential. 
While the density field describes where matter accumulates, the velocity field 
describes how matter moves in response to the evolving potential. The divergence 
of the velocity field, $\theta = \nabla \cdot \mathbf{v}$, plays a central role in 
structure formation, linking density growth to the temporal evolution of the 
gravitational field.

In the standard cosmological model, the velocity divergence evolves according to 
the linear growth rate and the continuity equation. In RDCC, the gravitational 
potential evolves differently. The relaxation mechanism suppresses the decay of 
small-scale modes and enhances the coherence of large-scale modes. This modifies 
the growth of the velocity divergence, the amplitude of bulk flows, and the 
coupling between density and velocity fields. The velocity divergence becomes a 
direct tracer of the relaxation scale $k_{\mathrm{rel}}$.
% ---------------------------------------------------------
\section{The RDCC Potential and the Continuity Equation}

The continuity equation relates the density contrast $\delta$ to the velocity 
divergence:
\begin{equation}
    \frac{\partial \delta}{\partial a}
    + \frac{1}{aH}\,\theta = 0.
\end{equation}

The velocity divergence is governed by the Euler equation,
\begin{equation}
    \frac{\partial \theta}{\partial a}
    + \frac{\theta}{a}
    = -\frac{k^{2}}{aH}\,\Phi_{\mathrm{RDCC}}.
\end{equation}

In RDCC, the gravitational potential takes the form
\begin{equation}
    \Phi_{\mathrm{RDCC}}(k,a)
    = \Phi(k,a)
      \left[
        1 - \chi_{\Phi}
        \left(\frac{k}{k_{\mathrm{rel}}}\right)^{\alpha}
      \right].
\end{equation}

This modification suppresses the contribution of small-scale modes to the velocity 
divergence while preserving the coherence of large-scale modes.
% ---------------------------------------------------------
\section{Growth of the Velocity Divergence Field}

The growth of the velocity divergence field is determined by the growth rate,
\begin{equation}
    f(a) = \frac{d\ln D}{d\ln a},
\end{equation}
where $D(a)$ is the linear growth factor.

In RDCC, the growth rate becomes scale-dependent:
\begin{equation}
    f_{\mathrm{RDCC}}(k,a)
    = f(a)
      \left[
        1 - \chi_{f}
        \left(\frac{k}{k_{\mathrm{rel}}}\right)^{\alpha}
      \right].
\end{equation}

Small-scale modes grow more slowly, while large-scale modes retain their coherence. 
The velocity divergence field becomes smoother and more coherent on scales near 
$k_{\mathrm{rel}}$, reflecting the temporal structure of the RDCC gravitational 
field.
% ---------------------------------------------------------
\section{Bulk Flows and Large-Scale Coherence}

Bulk flows arise from the coherent motion of matter on large scales. In RDCC, the 
enhanced coherence of large-scale modes amplifies bulk flows, while the suppression 
of small-scale modes reduces small-scale velocity fluctuations.

The bulk flow velocity on scale $R$ is given by
\begin{equation}
    v_{\mathrm{bulk}}^{2}(R)
    = \int dk\,
      P_{\theta}(k)\,
      W^{2}(kR),
\end{equation}
where $P_{\theta}(k)$ is the velocity divergence power spectrum.

In RDCC, this spectrum becomes
\begin{equation}
    P_{\theta,\mathrm{RDCC}}(k)
    = P_{\theta}(k)
      \left[
        1 - \chi_{\theta}
        \left(\frac{k}{k_{\mathrm{rel}}}\right)^{\alpha}
      \right].
\end{equation}

This modification enhances bulk flows on large scales and suppresses small-scale 
velocity fluctuations, producing a distinctive RDCC signature.
% ---------------------------------------------------------
\section{The Dynamical Architecture of the Cosmic Web}

The cosmic web is shaped not only by the density field but also by the velocity 
field. Filaments form where matter flows along coherent velocity streams, while 
voids expand where the velocity divergence is positive.

In RDCC, the enhanced coherence of large-scale modes produces more ordered flows 
along filaments and smoother expansion in voids. The velocity divergence field 
becomes a dynamical map of the relaxation scale, with coherent structures emerging 
on scales near $k_{\mathrm{rel}}$.

This dynamical architecture provides a new perspective on the formation of the 
cosmic web in RDCC.
% ---------------------------------------------------------
\section{Conclusion}

The velocity divergence field in the Relaxation-Driven Cyclic Cosmology reflects 
the scale-dependent evolution of the gravitational potential. The relaxation 
mechanism modifies the growth of the velocity divergence, the amplitude of bulk 
flows, and the coupling between density and velocity fields. These effects produce 
a distinctive dynamical architecture in the cosmic web, with enhanced coherence on 
large scales and suppressed fluctuations on small scales. The present Companion 
establishes the theoretical foundation for future studies of cosmic flows, 
redshift-space distortions, and the dynamical evolution of large-scale structure 
in RDCC.

% ========================
% References
% ========================

\begin{thebibliography}{10}

\bibitem{Flagship} Lehmann, M. (2026). Relaxation-Driven Cyclic Cosmology (RDCC): A Minimal CPT-Symmetric Framework with a Single Parameter. Flagship Paper. Zenodo: \url{https://zenodo.org/records/19744958} (v43.0 or later).

\bibitem{Zenodo} The complete RDCC companion series (I--LVII) is available at the same Zenodo record above.

% Thematisch relevante Companions
\bibitem{CompXVI} Companion XVI: Boltzmann Code and Modified Hierarchy.
\bibitem{CompXXXI} Companion XXXI: RDCC Halo Model and Non-Linear Structure.
\bibitem{CompXXXII} Companion XXXII--XXXIX: Galaxy--Halo Connection, Cosmic Web, Lensing, RSD, ISW, tSZ, Rotation Curves, CGM.

% Wichtige externe Referenzen
\bibitem{Planck2020} Planck Collaboration (2020), Astron. Astrophys. 641, A6.
\bibitem{DESI2024} DESI Collaboration (2024/2025), arXiv: [aktuelle $f\sigma_8$/BAO-Paper].
\bibitem{JWST2024} Maiolino et al. (2024), Nature (JWST high-$z$ galaxies/quasars).
\bibitem{Kaiser1987} Kaiser, N. (1987), Mon. Not. R. Astron. Soc. 227, 1 (RSD).
\bibitem{Bartelmann2001} Bartelmann, M. \& Schneider, P. (2001), Phys. Rep. 340, 291 (Weak Lensing Review).
\bibitem{HuOkamoto2002} Hu, W. \& Okamoto, T. (2002), Astrophys. J. 574, 566 (CMB Lensing).
\bibitem{LewisChallinor2006} Lewis, A. \& Challinor, A. (2006), Phys. Rep. 429, 1 (CMB Lensing Review).
\bibitem{NFW1997} Navarro, J. F., Frenk, C. S. \& White, S. D. M. (1997), Astrophys. J. 490, 493 (Halo Profiles).

\end{thebibliography}

\end{document}
